\section{Overview}
\label{sec:overview}

Fig.~\ref{fig:overview} summarizes our pipeline.
We first run a native-fine teacher in its approximately $2.399\,\mu$m CT frame.
Its probabilities, support, and medial crest are projected into the
$9.362\,\mu$m grid as soft training fields (Sec.~\ref{sec:crossres}).
We initialize every candidate from the exact released M7 checkpoint and train
all parameters under a small global trust region. The output is one ordinary M7 network.
At inference we use only its raw class-1 probability; there is no teacher and no
M7 probability blend.

At the selected threshold, an optional C postprocess operates on each axial
slice (Sec.~\ref{sec:fitter}). It proposes Bezier completions between skeleton endpoints, but paints a join
only when local curve geometry, radial displacement, neighboring slices, and
connection-track support agree. Segmentation training and curve completion are measured separately so an
improvement in one stage cannot hide a failure in the other.
All segmentation qualification numbers in Sec.~\ref{sec:results} use the raw
student without morphology or line fitting.

\begin{figure*}[t]
  \centering
  \figureasset{figure_1_method_overview}{0.96\linewidth}{2.05in}
  \caption{Method overview. The final figure will separate training-only
  components (fine teacher, projection, auxiliary losses, and frozen M7 anchor)
  from deployed components (raw student and optional 2D fitter). It will also
  show that the parameter trust-region projection acts after every optimizer
  update.}
  \Description{Reserved schematic of training-only and deployed pipeline components.}
  \label{fig:overview}
\end{figure*}

Two engineering details prevent accidental leakage or model substitution. First, the dynamic connectivity atlas is built only from boxes in the training
manifest; held-out morphology is used only for evaluation. Second, the M7 initializer is accepted only if both its byte size and SHA-256
match the released artifact. A new candidate may resume its own interrupted optimizer state, but may not use
an earlier student as learned initialization.
